\begin{mistake}{2026-06-05}{03}

\begin{problemcontent}
确定常数 \(a, b\)，使 \(x \to 0\) 时 \(f(x) = e^x - \frac{1+ax}{1+bx}\) 为 \(x\) 的三阶无穷小。
\end{problemcontent}

\begin{solutioncontent}
将原式通分：
\[
f(x) = e^x - \frac{1+ax}{1+bx} = \frac{e^x(1+bx) - (1+ax)}{1+bx}
\]
当 \(x \to 0\) 时，分母 \(\lim_{x \to 0} (1+bx) = 1 \neq 0\)。
因此 \(f(x)\) 的无穷小阶数与其分子相同，分子 \(g(x) = e^x(1+bx) - (1+ax)\) 必须是 \(x\) 的三阶无穷小。

对 \(e^x\) 使用麦克劳林公式展开至三阶：
\[
\begin{aligned}
g(x) &= \left( 1 + x + \frac{1}{2}x^2 + \frac{1}{6}x^3 + o(x^3) \right)(1+bx) - (1+ax) \\
&= \left( 1 + x + \frac{1}{2}x^2 + \frac{1}{6}x^3 \right) + bx\left( 1 + x + \frac{1}{2}x^2 \right) - 1 - ax + o(x^3) \\
&= (1 - 1) + (1 + b - a)x + \left( \frac{1}{2} + b \right)x^2 + \left( \frac{1}{6} + \frac{b}{2} \right)x^3 + o(x^3)
\end{aligned}
\]
由于 \(g(x)\) 是 \(x\) 的三阶无穷小，所以其一次项和二次项系数必须为 \(0\)：
\[
\begin{cases}
1 + b - a = 0 \\
\frac{1}{2} + b = 0 
\end{cases}
\]
解得：
\[
b = -\frac{1}{2}, \quad a = \frac{1}{2}
\]
此时三次项系数 \(\frac{1}{6} + \frac{b}{2} = -\frac{1}{12} \neq 0\)，符合题意。
\end{solutioncontent}

\mistakeknowledge{
\begin{itemize}
  \item \textbf{核心方法}：遇到分式加减求阶数，应先通分，若分母极限不为0，则阶数由分子决定。
  \item \textbf{易错点}：未通分前，不可对分母中的局部式子（如 \(1+bx\)）求极限强行代入，这会丢失高阶项导致方程缺失。
  \item \textbf{\(n\)阶无穷小条件}：泰勒展开后，\(0\) 到 \(n-1\) 次项系数全为 \(0\)，且 \(n\) 次项系数不为 \(0\)。
\end{itemize}}

\end{mistake}
